\section{Conclusion}\label{sec:short_conclusion}

We asked whether accounting for uncertainty about volatility changes
WTI average-price option values as much as choosing a different
volatility input. These are separate decisions. Posterior integration
changes the point price according to the spread of the volatility
posterior and the curvature of the pricing function. The spread of
the resulting model prices is a different measure of uncertainty.
Our synthetic results illustrate this distinction, while the
October-2026-expiry sample shows that posterior integration and
pricing at the posterior mean give almost identical point values.

In the original date-ordered study, using earlier option-implied
volatility yields smaller errors relative to settlement-price
proxies than the tested historical-return inputs. For the same
contracts, simply carrying forward the last available implied
volatility also has lower pooled error than fitting a new smile.
That finding is based on 15 valuation dates and is not a general
ranking of forecasting methods. A separate, longer Databento
exercise supports the importance of the volatility input under
a different pricing approximation. In that exercise, Bayesian
updating of earlier implied volatility does not reliably improve
on simple persistence. Bayesian intervals for settlement-price
errors also fall short of their intended coverage in parts
of early 2026.

For valuation practice, the main lesson is to compare a
posterior-integrated price with a simple plug-in price and to
test historical-volatility estimates against available earlier
option-implied volatilities. Uncertainty across model prices
should be reported separately from intervals for observed
pricing errors. More elaborate inference is useful only if
it improves the quantity being evaluated.

These conclusions concern conditional comparisons with
end-of-day settlement-price proxies, not executable quotes
or prospective trading profits. Some marks may be generated
from exchange valuation models; the longer exercise also uses
other same-date option prices as inputs. The original study
has relatively few valuation dates and many zero-volume
observations, and the historical publication times of vendor
records remain unverified. Better time-stamped data, more
actively traded contracts, and tests with richer futures
dynamics are needed before drawing stronger conclusions
about hedging or the volatility risk premium.
